\section{The new country, and what it turned out to measure}
\label{sec:china}

The study was complete when a constraint was imposed from outside it: add a
country. China was chosen because six provinces spanning Hainan to Heilongjiang
cover a wider climate range than the entire European panel and do so in a
developing economy, which is the comparison India actually needs. The protocol
adaptation it forced is stated in Table~\ref{tab:china} and is a single change:
the source is one calendar year, so the fifteen-month training block does not
exist and the calendar \emph{position} of the split is preserved instead ---
validate on April and May, test on June, as in every other row --- leaving three
months of training rather than fifteen. That biases against us, which is the safe
direction to err.

\begin{table}[t]
\centering
\small
\begin{tabular}{lrrlrrr}
\toprule
Province & Days & Shapes & Days/shape & Numbers & Values & Max err.\ \% \\
\midrule
Hainan (Haikou) & 365 & \textbf{2} & 115/250 & 778 & 8,760 & 0.000000 \\
Guangdong (Guangzhou) & 365 & \textbf{2} & 115/250 & 778 & 8,760 & 0.000000 \\
Shanghai & 365 & \textbf{2} & 115/250 & 778 & 8,760 & 0.000000 \\
Beijing & 365 & \textbf{2} & 115/250 & 778 & 8,760 & 0.000000 \\
Yunnan (Kunming) & 365 & \textbf{2} & 115/250 & 778 & 8,760 & 0.000000 \\
Heilongjiang (Harbin) & 365 & \textbf{2} & 115/250 & 778 & 8,760 & 0.000000 \\
\bottomrule
\end{tabular}
\caption{Provenance audit of the China arm, run against the data rather than taken from its documentation. Every calendar day's 24-hour profile is min--max normalised and the distinct shapes counted. \emph{Shapes} is how many survive; \emph{Numbers} is how many values are needed to rebuild the year from them; \emph{Values} is how many the year contains. A whole year of hourly provincial demand is 778 numbers --- 2 normalised day-shapes and one (min, max) pair per day --- reproducing all 8,760 values to within 2e-09 of their magnitude, which is the precision the source is published at rather than an approximation we introduced. Compression 11.26:1. The published method is exactly what the authors state it is. The consequence for this study is that forecast skill, a ratio against a seasonal-naive baseline, is uninformative on these series: the shape the baseline has to guess is constant, so both models reduce to predicting the same two administrative numbers per day.}
\label{tab:china-audit}
\end{table}


\paragraph{The audit.} This is the only tier-2 source whose documentation
describes a construction rather than a measurement. Wu and Kan take each day's
maximum and minimum from National Development and Reform Commission reporting,
trace a representative workday profile and a representative holiday profile off
published load curves with WebPlotDigitizer, and rescale one of those two
profiles to each day's envelope \citep{wu2023china}. Taking that on trust is
exactly what this project does not do elsewhere, so we checked it against the
data: min--max normalise every calendar day's 24-hour profile and count the
distinct shapes in a year.

The answer is two, in every province. Table~\ref{tab:china-audit} reports it:
$778$ numbers --- two normalised day-shapes and $365$ (min, max) pairs ---
reproduce all $8{,}760$ hourly values to within $2\times10^{-9}$ of their
magnitude, which is the precision the source is published at rather than an
approximation we introduced, at a compression of better than eleven to one. The
documentation is accurate.

\paragraph{What that disqualifies.} Skill in this study is a ratio against each
series' own seasonal-naive baseline. When every workday shares one shape,
seasonal naive recovers that shape exactly and its entire error is the daily
envelope --- and so is ours. Both models collapse onto predicting the same two
administrative numbers per day, the ratio between them collapses toward zero,
and the observed mean skill of $+0.040$ with two negative rows is what that
predicts rather than a finding about a forecaster. The China skill column is
reported for completeness and is not evidence either way about the method.

This costs us the arm's most quotable result. Guangdong is the closest structural
analogue to Delhi anywhere in the study --- cooling share $0.85$ against Delhi's
$0.89$, load--temperature correlation $+0.37$ against $+0.59$, a comparably peaky
profile --- and it returns $-0.314$, the second-worst skill in the benchmark.
Read naively that corroborates the Indian result on independent 2018 data from
another country and retires the objection that Delhi's weakness is an artefact of
a decade-old archive. It does not, because the number is a property of the
construction. We report the temptation as well as the retraction, because the
inference is one a reader would otherwise make on our behalf.

\paragraph{What survives.} Coverage is a property of interval width against
realised error and does not run through the baseline, so the tier-2 result can be
read on this arm --- and it replicates twice over. Under the benchmark's
calibration the six provinces average $0.824$ against $0.829$ on the metered
panel, on an error process unlike any other in the study, with the worst row
anywhere at $0.565$; under the adaptive step stated in the interval's units
they average $0.925$ against $0.914$, and the worst row, Heilongjiang, is
$0.937$ (Table~\ref{tab:aci-step}). Two populations sharing no country, no
year and no data-generating process landing within $0.005$ of each other under
the inert layer and within $0.011$ under the working one is stronger evidence
for both readings than a larger metered sample would have been: the failure is
universal, and so is the remedy. The null of Section~\ref{sec:axes} replicates
too. The wider coverage spread on this panel under the benchmark's calibration
($0.565$--$0.960$ against $0.762$--$0.911$) is consistent with the three-month
calibration history and is not separately identified.

\paragraph{What the exercise is worth.} A constraint imposed late produced no new
capability and one new measurement, which is a characterisation of the dataset
rather than of the method, and which removed a result we would have liked to
keep. That is the outcome we would want from an audit, and it is reported at the
same size as everything else.
